\chapter{The Claims of Community}

Michael Sandel begins this lecture by taking a step backward before moving forward. He does not begin with patriotism or loyalty or the moral claims of community. He first returns to the older contrast between Aristotle and Kant, because the communitarian challenge will make sense only if we see clearly what it is a challenge to. The argument therefore unfolds in stages: from law to freedom, from freedom to the self, from the self to the taxonomy of obligation, and only then to the test cases where loyalty and membership do real work.

\section{Kant's Reply to Aristotle}

Sandel opens with Kant's objection to Aristotle. Aristotle had argued that to think about the ideal constitution we must first think about the best way to live. Kant rejects that sequence. It is one thing, he says, for law to secure a fair framework of rights within which people pursue their own conceptions of the good. It is another thing altogether to build law itself on some favored conception of the good life. That second move, for Kant, carries the risk of coercion.

This gives us not only two views of politics but two views of freedom:
\begin{align}
\text{Aristotelian freedom} &= \text{living up to one's potential}, \\
\text{freedom} &\Longrightarrow \text{fit between persons and appropriate roles}, \\
\text{Kantian freedom} &= \text{autonomy}, \\
\text{autonomy} &= \text{acting according to a law I give myself}.
\end{align}

The same contrast appears at the level of law:
\begin{align}
\text{Aristotelian politics} &\Longrightarrow \text{law shapes character and cultivates virtue}, \\
\text{purpose of law / constitution} &= \text{fair framework of rights}.
\end{align}

Sandel is careful here. He does not merely report that Aristotle and Kant disagree. He shows where the disagreement lives. For Aristotle, politics aims at the good life because freedom itself is tied to realizing one's capacities and finding the role one is cut out for. For Kant, politics must not be organized around any such answer, because freedom is not role-fit but self-legislation.

\subsection{Question \& Answer}

\paragraph{Question.} What makes a person free: realizing a potential or giving oneself a law?

\paragraph{Answer.} Sandel wants us to feel that the two answers are genuinely opposed. Aristotle links freedom to fulfillment: to live freely is to live up to one's potential. Kant links freedom to autonomy: to live freely is to act according to a law one gives oneself. Once that difference is in place, the difference about law and justice follows naturally.

\section{The Unencumbered Self and Its Appeal}

Only after drawing the institutional contrast does Sandel widen the issue into a contrast about the person. Part of the appeal of Kant and Rawls, he says, lies in their picture of the self as free, independent, and capable of choosing its own ends. The liberal self is not simply a technical premise. It is morally attractive.

The lecture compresses that picture into a remarkably spare form:
\begin{align}
\text{self} &= \text{free and independent chooser of ends}, \\
\text{unchosen ties of history, tradition, inherited status} &\not\Longrightarrow \text{binding obligation}.
\end{align}

This is the liberating force of the liberal view. If we are free moral persons, then we are not bound simply by the accidents of birth. We are not morally constrained merely because we were born into a certain family, nation, religion, or tradition. We are, as Sandel puts it, the authors of the only obligations that constrain us.

That is why Sandel does not caricature the liberal position. He pauses over its attraction. The self appears here as sovereign over its ends, unbound by ties it has not chosen. If the communitarian critique is to carry weight, it must show not merely that people often feel attached to history and membership, but that this liberal self leaves out something real in our moral experience.

\subsection{Question \& Answer}

\paragraph{Question.} Why is the image of the free and independent chooser so powerful?

\paragraph{Answer.} Because it promises freedom from inherited rank and unchosen moral burdens. It says that history and status do not bind us until we endorse them. Sandel grants the force of that vision before criticizing it, and that order matters: communitarianism enters not as a denial of freedom, but as a complaint that this account of freedom is incomplete.

\section{MacIntyre and the Narrative Conception of the Self}

Sandel now turns to Alasdair MacIntyre, not merely as a critic of liberalism but as the source of a different picture of the self. Human beings, MacIntyre argues, are storytelling creatures. We do not answer the practical question from nowhere. We answer it from within a history.

That move is the hinge of the lecture:
\begin{align}
\text{What am I to do?} &\Longrightarrow \text{What story or stories do I find myself a part of?}, \\
\text{good for me} &= \text{good for someone who inhabits these roles}, \\
\text{past} &\Longrightarrow \text{debts}+\text{inheritances}+\text{expectations}+\text{obligations}, \\
\text{narrative self} &\approx \text{encumbered self}.
\end{align}

The point is not merely that we have biographies. It is that our biographies are morally relevant. I am someone's son or daughter, a citizen of this or that city, a member of this nation or that people. What is good for me cannot be abstracted from those roles, because the self that asks the question is already situated within them.

Sandel follows MacIntyre one step further. The narrative self is not only socially located; it is morally claimed. The history, tradition, and communities of which I am a part help constitute what I owe. That is why MacIntyre treats historical amnesia as a moral failure. To say, ``I never owned any slaves,'' or ``I was born after 1945,'' may register a biological fact, but it may also refuse a history that still shapes present responsibility.

\subsection{Question \& Answer}

\paragraph{Question.} Why can I not answer what I ought to do without asking what story I am part of?

\paragraph{Answer.} Because, on the narrative view, the self does not begin as an empty chooser. It begins already situated in roles, relationships, memories, and inheritances. The practical question therefore arrives already burdened with history. Sandel's claim is not that history decides everything, but that moral reflection cannot pretend to begin without it.

\section{From Two Kinds of Obligation to a Third}

At this point Sandel makes the lecture's cleanest formal move. He says that on the liberal conception moral and political obligations arise in one of two ways. First, there are natural duties: duties we owe human beings as such. Second, there are voluntary obligations: duties we incur by promise, agreement, contract, office, or some other act of consent.

For compactness we write
\begin{align}
D_N &:= \text{natural duties}, \\
D_V &:= \text{voluntary obligations}, \\
D_S &:= \text{obligations of solidarity, loyalty, or membership}.
\end{align}
Then Sandel's taxonomy becomes
\begin{align}
\text{natural duties} &= \text{duties owed to human beings as such}, \\
\text{voluntary obligations} &= \text{obligations arising from promise, deal, or contract}, \\
\text{liberal obligations} &= \text{natural duties}+\text{voluntary obligations}, \\
\text{communitarian obligations} &= \text{natural duties}+\text{voluntary obligations}+\text{solidarity / loyalty / membership}.
\end{align}

This is the formal spine of the lecture. The dispute is no longer simply over whether community matters. It is over whether the liberal two-term classification is exhaustive. The communitarian says no. There is a third category, and without it we cannot account for certain obligations we commonly recognize and even prize.

\paragraph{Worked derivation.} Sandel's reasoning can be written as a short chain:
\begin{enumerate}
\item If the self is fundamentally free and independent, then unchosen history does not bind simply as such.
\item If unchosen history does not bind, obligations arise either universally or voluntarily.
\item Hence the liberal scheme \(D_N + D_V\).
\item If, however, the self is partly constituted by inherited memberships, then some obligations may arise neither from universality nor from consent.
\item Hence the communitarian addition of \(D_S\).
\end{enumerate}

Sandel's caution is important. At this stage he has not yet proved that \(D_S\) is real. He has only put the issue in its cleanest form. The rest of the lecture will test whether this third term names a genuine moral phenomenon or merely bundles together sentiment, reciprocity, and prejudice.

\subsection{Question \& Answer}

\paragraph{Question.} Is there a third category of obligation beyond universal duty and consent?

\paragraph{Answer.} This is now the exact question. Sandel has moved from a dispute about the self to a dispute about the types of obligation we can recognize. The communitarian claim is not vague warmth toward community. It is the sharper claim that the liberal classification leaves out a distinct and indispensable kind of moral tie.

\section{Test Cases of Membership: Family, War, Airlift, Patriotism}

Sandel does not leave the taxonomy abstract. He now turns to examples, and the order matters. He begins with family, where the pressure on the liberal picture is strongest, and then moves outward to explicitly political cases.

First comes the parent-child case. If two children are drowning and only one can be saved, one your own and one a stranger's, must you flip a coin? Sandel asks the question in order to provoke the thought that strict impartiality may itself be morally obtuse. Then he turns the example around. Parents may have chosen to have children, but children do not choose their parents. If one aged parent is yours and one belongs to a stranger, why does it seem morally intelligible that your obligation falls more heavily on your own parent? That asymmetry is precisely the point. The obligation appears not to be traceable to consent.

The political examples then extend the same line of thought. A French resistance pilot refuses to bomb his own village, not because he doubts the justice of the cause, but because it would be a special moral crime for him to bomb his own people. The Israeli airlift of Ethiopian Jews sharpens the issue differently. Is the rescue of fellow Jews from famine a prejudice, or the recognition of a special obligation of solidarity? Sandel does not settle these cases yet. He uses them to make \(D_S\) morally vivid.

That sequence leads him to patriotism. Two towns face one another across the Rio Grande, Franklin, Texas and Franklin, Mexico. Why, if at all, should national boundaries matter morally? Why should needy people on one side of the river claim us more strongly than equally needy people on the other? Sandel formulates the communitarian answer this way:
\begin{equation}
\text{patriotism at least potentially a virtue}
\Longrightarrow
\text{expression of obligations of citizenship}.
\end{equation}

Only now does the lecture ask the room whether it is sympathetic to this third category of obligation. The conceptual claim has been translated into lived cases first.

\subsection{Question \& Answer}

\paragraph{Question.} Why should family or country bind us more strongly than strangers if consent did not create the tie?

\paragraph{Answer.} The communitarian answer is that these ties are not merely external relations; they help constitute the person who deliberates. Sandel is careful, though, not to let that answer settle too quickly. The point of the examples is to raise the question sharply enough that the objection can now be stated in earnest.

\section{First-Round Classroom Objections}

Once the cases are on the table, the lecture shifts into live argument. Sandel begins with the critics. The first objection is structural. If obligations of solidarity are real, then we inhabit many communities at once, and their claims can collide. Patrick's worry is not that loyalty never matters, but that a framework built on solidarity may have no clear way to rank competing loyalties.

That worry generates two opposed proposals, neither of which Sandel treats as decisive:
\begin{align}
\text{most universal community} &\stackrel{?}{>} \text{more particular communities}, \\
\text{family} &\stackrel{?}{>} \text{town} \stackrel{?}{>} \text{country} \stackrel{?}{>} \text{humanity}.
\end{align}

Nicola suggests that the most universal community, the human community, should take precedence. Patrick replies that this itself seems arbitrary, and proposes the opposite direction: perhaps the most immediate and particular community should come first. Sandel lets both proposals surface because together they reveal the difficulty. Scale alone does not yet give us a principle.

The next objection concerns patriotism more specifically. Elizabeth argues that citizenship is a constructed narrative, perhaps even a false one. A river or a national boundary may be a historical accident. Why should the lottery of birth generate morally weighty obligation? In reply, other students try to redescribe the relevant obligations in terms of kinship and reciprocity rather than sheer membership. The more closely we interact with people, the thought runs, the stronger the reciprocity; perhaps that is enough.

Rena then presses a different line. Some attachments, she suggests, may resemble school spirit or house spirit: emotionally real, but not yet morally weighty. Sandel answers by returning to the parent case. If the tie to an aged parent is more than sentiment, what exactly makes it more? Rena's answer pushes again toward benefit and repayment rather than solidarity as such.

\subsection{Question \& Answer}

\paragraph{Question.} If communities overlap, how are competing claims of loyalty to be ranked?

\paragraph{Answer.} The lecture's first answer is that neither the universal-first rule nor the particular-first rule is obviously satisfactory. Sandel uses this impasse to show that the communitarian view cannot simply appeal to membership and stop there. Once loyalties overlap, a further moral judgment is required.

\section{Recap Seam and the Strongest Objections Reframed}

At 00:24:26 the lecture begins again. Sandel explicitly says that he now wants to consider the strongest objections to obligations of solidarity and membership. This is not a repetition of the earlier discussion. It is a reset, and the reset matters. The first round had exposed the pressure points; this second pass gathers them into a more systematic form.

The objections now come into focus as four:
\begin{enumerate}
\item obligations of membership may conflict;
\item the cases may show sentiment rather than obligation;
\item the apparent obligations may be reducible to consent or reciprocity;
\item communal obligation may amount to collective selfishness or prejudice.
\end{enumerate}

The liberal reduction is now stated as clearly as possible. Julia's Rawlsian point is that patriotism and familial love can be honored within a liberal framework as voluntary commitments. Rawls himself, Sandel notes, makes room for political obligation when someone knowingly undertakes an office or enlists in the military. But for citizens as such there is, strictly speaking, no political obligation absent some binding act of consent. The same structure then extends to family and country more generally:
\begin{equation}
\text{consent} \lor \text{reciprocity}
\stackrel{?}{\Longrightarrow}
\text{all particular obligations}.
\end{equation}

So the second-round liberal claim is stronger than the first. It does not deny that we may love our family or country. It says that all the moral force of these attachments can be absorbed by already familiar liberal categories. Particular loyalties are permitted if chosen, and permitted if they do not violate our natural duties to persons as persons. No third category is needed.

Sandel adds one more objection before moving on: perhaps obligations of membership are simply a kind of collective selfishness. Why should we honor them at all? Why are they not just prejudice under a warmer name?

\subsection{Question \& Answer}

\paragraph{Question.} Are solidarity and patriotism genuine moral obligations, or just emotion, consent, or reciprocity under another name?

\paragraph{Answer.} Sandel's answer here is still provisional. He does not yet refute the liberal objection. What he does instead is sharpen the test. If every compelling case of loyalty can be redescribed as sentiment, consent, reciprocity, or harmless preference, then \(D_S\) disappears. So the lecture now needs a case in which loyalty genuinely competes with universal moral claim.

\section{Patriotism, Hard Cases, and the Return to Justice}

The lecture now turns from the abstract taxonomy back to patriotism, but on more demanding terms. Sandel invites defenders of patriotism to state the strongest communitarian case they can. A.J. argues that country, like family, forms identity. One can owe more to one's country than to humanity in general without that being mere prejudice, just as one may owe more to one's parents than to someone else's. Rahul develops the point in a civic direction. Patriotism, he says, can sustain debate, criticism, and shared civic responsibility; one can love one's country and oppose its present policies.

Sandel lets the critics answer. If patriotism is reduced to critical attachment to just principles, they say, then what remains of loyalty as such? Julia presses exactly this point. The defenders of community need a case where loyalty does more than echo universal justice. It must actually compete with it. That is the harder test.

The lecture then produces that test. Dan says he would not report a cheating roommate. Billy Bulger refuses to help authorities catch his fugitive brother. Robert E. Lee refuses to fight against Virginia, even though he opposes secession and even though the cause opposed includes slavery. Sandel chooses these cases because they are morally dangerous. If loyalty still seems morally weighty here, then it cannot be dismissed as a mere sentimental afterglow.

\paragraph{Worked test.} The hard cases line up in a clear progression:
\begin{enumerate}
\item Dan's dilemma: loyalty to roommate competes with fairness and truth-telling.
\item Billy Bulger's dilemma: brotherly loyalty competes with civic and legal responsibility.
\item Robert E. Lee's dilemma: loyalty to Virginia competes with Union and, ultimately, with a just cause against slavery.
\end{enumerate}

Sandel's point is precise. The communitarian need not say that loyalty rightly wins in every such case. The claim is narrower and more unsettling: there may be something morally intelligible, perhaps even admirable, in the pull of loyalty even when the ultimate act is gravely contestable. If that is so, then loyalty is not merely sentiment.

The critics answer as Sandel wants them to answer. They say these cases reveal multiple, competing communities and no single communitarian principle for ranking them. One student proposes the most local community as morally weightiest; another rejects scale as the decisive factor altogether. Sandel then performs his final turn. If the communitarian challenge is right, and if justice cannot be detached from conceptions of the good embodied in shared forms of life, what follows for justice itself?

This is where Michael Walzer enters. Sandel introduces him as the strongest articulation of the communitarian implication:
\begin{align}
\text{justice} &= \text{relative to social meanings}, \\
\text{just society} &= \text{faithfulness to the shared understandings of its members}.
\end{align}

Now the issue becomes larger than patriotism. Suppose the priority of the right over the good cannot be sustained. Suppose justice and rights are unavoidably bound up with conceptions of the good as lived in particular communities. Must justice then collapse into convention?
\begin{align}
\text{justice tied to shared meanings} &\Longrightarrow \text{conventionalism?} \\
\text{tradition} + \text{situated identity} &\not\Longrightarrow \text{justice}.
\end{align}

Sandel does not settle the matter. Instead he plays the clip from \emph{Eyes on the Prize}. Segregationist southerners invoke precisely the language that communitarianism seems to dignify: tradition, situated identity, inherited ways of life, shared understandings. That is why the lecture ends where it does. The question is no longer whether community matters. It is whether justice can be tied to shared meanings without licensing unjust traditions.

\subsection{Question \& Answer}

\paragraph{Question.} If justice follows shared social meanings, what prevents convention from ratifying injustice?

\paragraph{Answer.} This is the unresolved question on which the lecture closes. Walzer shows how justice might be socially embedded rather than abstractly detached. The segregation example shows the danger of that move. Sandel leaves us with the problem, not the solution: can shared meanings be morally authoritative without becoming morally tyrannical?

\section{Summary}

The lecture begins with Aristotle and Kant, but only in order to reach a more exact dispute about the self and about obligation. From the liberal self comes a two-part moral taxonomy, \(D_N + D_V\). From the narrative or encumbered self comes the communitarian claim that a third category, \(D_S\), must be added if we are to make sense of loyalty, solidarity, and membership. Sandel then tests that claim through family cases, wartime and rescue cases, patriotism, and finally the hard dilemmas of Dan, Billy Bulger, and Robert E. Lee. The last step is the widest one: if justice is bound to social meanings, can it survive bad traditions? The lecture ends, deliberately, with that question still open.